\begin{center}
\textsc{Gauss} to \textsc{Bessel}.\quad Göttingen, 15~November 1822.
\end{center}

\begingroup
\makeatletter
\renewcommand{\@makefnmark}{\hbox{\normalfont\textsuperscript{*})}}
\long\def\@makefntext#1{\parindent1.5em\noindent\textsuperscript{*})\ #1}
\makeatother
\par\noindent\hspace*{\parindent}.\ .\ .\ .\ .\ .\quad I already knew of the extraordinary difficulties of carrying a network of triangles through the Lüneburg Heath from \textsc{Epailly}'s report. He declares it outright impossible, and in order to connect the southern part of
\par\clearpage
\noindent352\hfill{\small\textbf{CORRESPONDENCE WITH BESSEL.}}\par\medskip
\noindent Hanover with Hamburg, he carried his triangles via Bremen, down the Weser and back up the Elbe.
And yet he had great advantages over me. He always observed from the tops of his signal towers, where he could far more easily obtain a wider view than I could, having measured everywhere at ground level. In cutting through the woods he had far less need to consider preserving the trees. He had free transport, free labor everywhere, etc., whereas I have to pay for everything without exception in cash, and often see the costs go very, very far beyond what I had expected. He had an entire brigade of engineers to assist him, etc.
These considerations made me undertake my first reconnaissance trip with some apprehension (toward the end of April).
At first things actually went better than I had expected. I found that Garssen and Falkenberg could be connected directly with Deister and Lichtenberg, but in investigating how the triangles could be carried farther north from those two points as far as Lüneburg, I found the terrain so unyielding that I several times doubted whether a successful outcome was possible and feared that I would have to give up the whole undertaking.
The land flat everywhere, no commanding points, woods everywhere, partly in large forests, such as the Hassel, the Lüsing, the Becklinger Holz, etc., partly in countless smaller enclosed stands of trees interposed like squares on a chessboard.
My assistant's attempts to find anything usable on the western side were quite unsuccessful. I myself finally succeeded after the most arduous attempts in establishing two triangles, 9.~10.~12 [Falkenberg-Hauselberg-Wulfsode] and 10.~12.~13 [Hauselberg-Wulfsode-Wilsede], in what one might call the heart of the heath. At the same time, however, I became convinced that \textls[120]{this} sort of work would soon prove too much for me, so I put off any further search for a way to link these triangles to those farther south and to continue them farther north until a later time, when I would have more assistants (I had taken only one officer on the first trip) and all my instruments with me.
I therefore began the actual measurements at Lichtenberg in mid-June, and my hopes have since been fairly well fulfilled.
Since it was found quite impossible to connect Garssen with Hauselberg by cutting a sight line through the woods, a point still farther north, 14 [Scharnhorst],
\par\clearpage
\noindent{\small\textbf{ON THE HANOVER TRIANGULATION.}}\hfill353\par\medskip
\noindent was found after a long search and offered more hope.
Later, however, this proved absolutely impossible, since the forest between 10 [Hauselberg] and 14 [Scharnhorst] lay on ground that was too high, or rather not low enough, and I had to count myself fortunate to find yet another point, Breithorn, with which I finally succeeded.
Meanwhile, the line 9.~13 [Falkenberg-Wilsede] had also been opened by a very substantial clearing, and thus 10 [Hauselberg] could have been omitted altogether if 11 [Breithorn] had been known earlier.
However, I retained the former point, and I consider it an exceedingly valuable advantage that my system includes three quadrilaterals in which all six directions were actually measured both ways.
After all the angles in each triangle had first been duly adjusted to the proper sum, only very small modifications, mostly below \(0^{\prime\prime}{,}1\), were needed to bring these three quadrilaterals into full agreement.
It would be desirable to have \textls[120]{such} checks for every survey.
There are surveys in which the sums of the three angles agree admirably everywhere, yet such a check shows that some angles are certainly wrong by \(2^{\prime\prime}\) to \(3^{\prime\prime}\).
Indeed, the check by means of the sum of the angles is within everyone's reach. The check by means of diagonals is less so, however easy it may be for a mathematician, and one cannot help suspecting that the former check may sometimes have served to \textls[120]{select} the observations to some extent, if not to falsify them (one notices a tendency toward this even in \textsc{Delambre}).
Fortunately, the Timpenberg was also found north of Wulfsode, and can be connected directly with Hamburg. The attempt to cut a sight line from Timpenberg to Lüneburg failed because the ground in between did not dip sufficiently.
The directions from 10 to 11 and from 11 to 14 required large clearings.
I always made it a rule to keep the calculations abreast of all the measurements as I obtained them (down to the very last line). Only in this way was it possible to carry out all the clearings with the utmost precision, so that not a single tree was felled needlessly, or to establish as early as possible that a clearing was impossible.
For example, I knew in advance the exact angle of depression at which 14 [Scharnhorst] would have to appear at 10 [Hauselberg] or 20 [Lüneburg at 19 [Timpenberg]. In the first
\par\smallskip\noindent{\footnotesize IX.}\hfill{\footnotesize45}\par
\clearpage
\noindent354\hfill{\small\textbf{CORRESPONDENCE WITH BESSEL.}}\par\medskip
\noindent case the ground was already too high in front of the woods, in the second after a narrow gap had been cut for 2000 paces.
From Wilsede I have fixed another point [Nindorf], 3000 meters north of Timpenberg, which has direct lines of sight to 20 [Lüneburg], 21 [Lauenburg], 23 [Hamburg] and can probably be connected with 19 [Timpenberg], but the entire stretch to it must be cleared.
At Timpenberg I have so far made only provisional measurements of the angles between 12 [Wulfsode], 13 [Wilsede], 23 [Hamburg]. All the points farther south have been completed, as has Wilsede.
Scharnhorst was the last place I visited. This point could be connected directly with Lichtenberg and Deister by cutting a few easy sight lines, which would make Garssen dispensable, but it would have been impossible to ascertain the usefulness of that point at the outset.
This point, like several others, is quite an inconspicuous spot, which one would never suppose to be of such value.
In many directions, for example from Lichtenberg to Garssen and from Lichtenberg to Falkenberg, the line passes so close above the intervening obstacles that with ordinary refraction [the target point] rises only \textls[120]{a few arcseconds} above them and only occasionally reached \(30^{\prime\prime}\) to \(40^{\prime\prime}\). From Falkenberg to Wulfsode and vice versa, with weaker refraction \textls[120]{no} light came through, and the light always rose into view only in the later hours of the afternoon.
From Lichtenberg itself I seldom saw anything on the Falkenberg and vice versa. The heliotrope light almost always hung in the open sky (distance 85542 meters). Despite the great distance, the measurements agree just as well, or rather almost better, than at quite short distances of 10000 meters, where the heliotrope light is still almost too strong even though the mirror was covered except for an opening of a few square millimeters.
To save time, even at distances of 3 to 4 to 5 miles, I often sighted on my stone pedestals themselves, which were \(3\frac{1}{2}\) feet high.
This made the measurements somewhat less accurate here and there than if I had used heliotrope light alone, but in that case I would not have got nearly so far this year.
The greatest error in the sum of the three angles this year was \(1^{\prime\prime}{,}76\), also the greatest \textls[120]{overall}\footnote{In 19 triangles.} after the one
\par\clearpage
\noindent{\small\textbf{ON THE HANOVER TRIANGULATION.}}\hfill355\par\medskip
\noindent in triangle 4.~5.~6 [Hohehagen-Hils-Brocken], where it amounts to \(3^{\prime\prime}{,}7\) and is chiefly attributable to the difficulty of sighting on the Brockenhaus tower, which in sunshine can never be bisected properly owing to its phase.
I would gladly measure the angles of this triangle again if time allowed.
This year I often had three heliotropes in operation at once, one of them tended by my son. It is indeed a magnificent sight.
The notice on heliotropes given by a certain \textsc{Schubach} in the [Astronomischen] Jahrbuch 1825 is based on an error and has nothing in common with my heliotropes. They are ingenious instruments whose design cost me \textls[120]{a very great deal} of trouble at first, but now, I believe, leave nothing to be desired.
I have had specimens of both designs, which are quite different from one another, made here for General \textsc{Müffling}. \textsc{Gerling} has also received two. Perhaps I will soon be able to give illustrations of them in \textsc{Schumacher}'s Astronomischen Nachrichten.
For signaling, and to show the heliotropes at the opposite stations the direction initially when the distances were great, I often used a large mirror one foot square, which was steered in yet another way.
According to my assistants' description, the sight of it was most magnificent. As far as four miles away it sometimes \textls[120]{hurt} the naked eye to look at it for long. --- But enough of this for now.

\par\noindent\hspace*{\parindent}I was pleased to see your way of treating geodetic observations in \textsc{Schumacher}'s Astronomischen Nachrichten.
You know that this subject occupied me many years ago.
Since it is hard to add anything to your work, if our paths had met, that publication would have made my own work superfluous.
But the way I have treated this subject differs from yours through and through, and so in publishing my measurements in the future I will also develop my theoretical work in detail.
I hope to be able to offer some unexpected things in it.
But these investigations are connected with a rich, almost inexhaustibly rich field, and on this as on so many other occasions I often feel with deep regret how my external circumstances prevent me from undertaking theoretical work of broad scope.
For such work to flourish fully, one must be able to devote oneself entirely to it and not
\par\smallskip\noindent\hfill{\footnotesize45\textsuperscript{*}}\par
\clearpage
\noindent356\hfill{\small\textbf{CORRESPONDENCE WITH BESSEL.}}\par\medskip
\noindent be hindered every hour by such disparate tasks as giving lectures, all the petty details of observing and computing the observations, etc.\ etc.\ .\ .\ .\ .\ .\ .

\par\noindent\hspace*{\parindent}The exertions in the hot summer have often been extremely taxing for me, at times so much so that I thought I would succumb to them.
Another great hardship in working on the desolate Lüneburg Heath is that often one can find only poor lodgings, and even those only miles from the place of work.
In cool weather, which suits my constitution better, I was generally in tolerably good health, and at present I cannot complain about how I feel.\ .\ .\ .\ .\ .\ .

\par\noindent\hspace*{\parindent}You can see from the above report what is still needed to complete my triangulation, but I cannot yet make a definite plan for next summer. Much will also depend on \textsc{Schumacher}'s cooperation.
It would be \textls[120]{possible} for me, if I confine myself to what is indispensable, to finish the triangles in April and May and then return here, perhaps to measure with the zenith sector here in July and at a point farther north (Celle, Harburg, Hamburg?) in August.
But if I do not need to spare expense, I may devote the whole summer to the triangulation to bring everything as close to perfection as possible.

\par\noindent\hspace*{\parindent}A few days ago I received a universal instrument from Munich which I had ordered two years ago chiefly for measuring azimuths.
Meanwhile, according to my provisional reduction, the azimuth derived from my transit instrument and carried over to Hamburg seems to differ by less than \(1^{\prime\prime}{,}5\) from the one \textsc{Schumacher} reported to me (he has also made a provisional measurement of the angle between Lüneburg and Wilsede, from where I had heliotrope light sent to him), and an azimuth measurement on my part at other points thus seems fairly superfluous.
Whether I will see reason to measure the altitudes of the pole at several points with the universal instrument will become apparent once I have used it for a while.
It is well made, but seems to me rather awkward to use, particularly for a nearsighted eye, which cannot get close enough to read the setting circles.
In any case, doing much observing with it must be very tiring.
All my hori-
\par\clearpage
\noindent{\small\textbf{ON THE HANOVER TRIANGULATION.}}\hfill357\par\medskip
\noindent zontal angles were measured with a 12-inch theodolite, which is an excellent instrument.
Last winter I made only a few measurements with the zenith sector, more or less by way of trial. It is far inferior to \textsc{Reichenbach}'s meridian circle in optical power and ease of use.

\par\noindent\hspace*{\parindent}My arc measurement as such can scarcely yield a very important result on its own, but I do not know when \textsc{Schumacher}'s may be expected to be complete.
As far as I know, he has carried his triangles only about one third of the total length so far, and the zenith-sector measurements at Skagen were made only by his assistant \textsc{Caroc}.
I do not yet know the exact length of his baseline (about 6000 meters). He meant to connect it to Hamburg only this autumn, but in his last letter he makes no mention of it.
I could, of course, measure a baseline myself, and the line from Breithorn to Scharnhorst seems to present no insurmountable obstacles along its entire length (11220 meters).
But even apart from the great expense, I confess I shrink from such exceedingly tedious work.
I have always been able to find something of interest in my trigonometrical measurements, since reducing them each day always provided some amusement\footnote{My practice in observing was always to measure whatever could be measured well at the time, regardless of whether it was an angle directly belonging to a triangle. At some stations I used a few auxiliary points so as not to be idle even when only one heliotrope was shining. The French have done something similar in Spain, as I see, but adjusted the measurements quite incorrectly.}.
In addition, whenever the opportunity arose I took intersecting sights on all visible objects to be of use to the geography of the region, and I must say that I grew so fond of this business with its daily adjustments that noticing a new church tower and working out and computing its position probably gave me just as much pleasure as observing a new celestial object.
(In the eyes of God it probably comes to the same thing in the end whether we have fixed the position of a church tower to within a foot or that of a star to within an arcsecond.)
But in measuring a baseline, once it is under way, I see absolutely nothing to engage or entertain the mind, and in strenuous daily work that may last two months one must sustain oneself solely with the thought
\par\clearpage
\noindent358\hfill{\small\textbf{CORRESPONDENCE WITH BESSEL.}}\par\medskip
\noindent that it simply \textls[120]{must} be done sometime in order to have a number in the end.

\par\noindent\hspace*{\parindent}The great accuracy of measuring horizontal angles by heliotrope light over the most immense distances, and the accuracy with which absolute azimuths can be measured with universal instruments, lead me to believe that it would now be possible to carry out a arc measurement in longitude with great advantage in suitable terrain by basing the difference in longitude on the convergence of the meridians instead of on time determination.
From the Pyrenees one can see Mont Ventoux, from there the Alps, from the Alps as far as Styria, from there as far as Hungary, etc.
Admittedly the accuracy here does not increase as the total length of the arc, as it does with a meridian arc (the geodetic part proper can be regarded as error-free in that case). Thus, if one imagines sections of roughly equal size, the accuracy increases as the number of sections for a degree of latitude, and likewise for a degree of longitude if the difference in longitude of the endpoints is observed astronomically, for instance by occultations of stars, but with the convergence of the meridians it increases only as the square root of the number of sections (approximately). Still, I believe that with \textls[120]{very large} sections this procedure gives greater accuracy than can be obtained, for example, by seeking the difference in longitude from gunpowder signals by means of \textls[120]{time determination}.
This procedure would be most advantageous in northern regions if there are very long lines of sight there.
Might such lines of sight be found in Norway and Sweden, or in some parts of Russia?

\par\noindent\hspace*{\parindent}.\ .\ .\ .\ .\ .\quad I am still at work on the adjustment, which I unfortunately did once entirely in vain, as an incorrect table from one station had been entered through an oversight.
In my treatment, every angle measured, for example, at Wilsede, in a sense affects all the others as far as Göttingen.
I will give you some results in the future.
According to a provisional calculation, taking Göttingen to be at \(51^{0}\,31^{\prime}\,48^{\prime\prime}{,}7\), Hamburg falls at \(53^{0}\,33^{\prime}\,1^{\prime\prime}{,}76\), and \(0^{0}\,2^{\prime}\,2^{\prime\prime}{,}97\) east of Göttingen, with the absolute value provisionally based on \textsc{Zach}'s baseline.
\par\noindent\hspace*{\parindent}.\ .\ .\ .\ .\ .
\endgroup
